\documentclass[10pt]{article}
\usepackage[T1]{fontenc}
\usepackage{lmodern,amsmath,amssymb,amsthm}
\usepackage[margin=1in]{geometry}
\usepackage[hidelinks]{hyperref}
\newtheorem{theorem}{Theorem}
\newcommand{\Z}{\mathbb Z}
\setlength{\parskip}{5pt}
\setlength{\parindent}{0pt}
\title{Rigidity and Inverse Problems\\for Restricted Lattice-Path Enumeration}
\author{Nicholas Ham}
\date{Mathematical briefing for possible discussion with James East\\16 October 2026 --- draft, not sent}
\begin{document}
\maketitle
\section*{1. From our earlier work to a small set of questions}
This revisits my 2019 conjectures using the distinction between words, endpoints and constrained transitions in our earlier paper [1]. Its first version already contains the factorial examples (Example 3.14); those are background. The present question is what regular endpoint counts reveal about permitted jumps. The attached draft gives proofs. I would value your view on whether the statements are familiar and whether any suggest worthwhile development; their originality and publication value remain unresolved.

Paths start at $(0,0)$, with empty weight one. A jump $(1,d)$ has time-independent weight $f(d)$; all vertices at layer $n$ must lie in a finite integer interval $I_n$. Thus
\[
p_n(y)=\sum_{z\in I_{n-1}}f(y-z)p_{n-1}(z).
\]
Indicator weights represent step sets, not multisets. A row is $q$-uniform if values agree at congruent admitted heights. Preservation of \emph{every} residue-constant input is stronger than uniformity of this one actual trajectory.

\textbf{Finite structural rigidity.} For $I_n=\{0,\ldots,n\}$ and $N\ge q$, every transition through $N$ preserves residue-constant inputs iff $f$ is $q$-periodic on $W_N=\{1-N,\ldots,N\}$. This holds over a commutative ring. At all horizons, indicator kernels are exactly unions of residue classes. The first informative horizon is $q$; earlier target residue cells are singletons.

For a source residue cell with extreme heights $a_r,b_r$, its predecessor sum satisfies
\[
C_r(y+q)-C_r(y)=f(y-a_r+q)-f(y-b_r).
\]
Comparing consecutive triangular layers isolates individual weights. For arbitrary intervals, join the two differences in each displayed equality. Preservation is exactly constancy on the resulting components; geometry forces periodicity iff each observable residue fiber is connected. The component formulation is an equality-constraint description. The invariant-space principle is classical equitable-partition theory [2]; the exact growing-window converse is the point requiring recognition.

\textbf{Widening funnels.} For any integer $c$ and $0=s_0\le\cdots\le s_N$, layers $I_n=[cn,cn+s_n]$ preserve constant inputs iff $f$ is constant on $[c-s_{N-1},c+s_N]$. Widths may jump arbitrarily. After shearing, a transition from width $M$ to $K$ gives edges of length $M+1$ on $[-M,K]$. The old connected window meets every class modulo $M+1$, connecting the new window. General-$q$ funnel rigidity is not claimed.
\newpage
\section*{2. Actual finite-range trajectories: obstruction and sharpness}
\begin{theorem}[Mass obstruction]
Let $f\ge0$ have finite nonempty support with extremes $m,M$. For $I_n=[A_n,B_n]$ put $J_n=I_n\cap[A_{n-1}+M,B_{n-1}+m]$. Suppose rows $b,\ldots,h$ are $q$-uniform, each layer has at least $q$ vertices, every intervening core $J_n$ has $q$ consecutive vertices, and
\[
\lfloor|I_h|/q\rfloor>\lceil|I_b|/q\rceil.
\]
Then row $b$ is zero.
\end{theorem}
\textbf{Short proof.} Let $c_n(r)$ be the residue values, $C_n=\sum_r c_n(r)$, $P_n=\sum_y p_n(y)$ and $K=\sum_d f(d)>0$. Core targets give $C_n=KC_{n-1}$; total outgoing mass gives $P_n\le KP_{n-1}$. Hence
\[
\lfloor|I_h|/q\rfloor C_h\le P_h\le K^{h-b}P_b
\le K^{h-b}\lceil|I_b|/q\rceil C_b.
\]
For $C_b>0$, substitution contradicts the strict gap. Nonnegativity handles $C_b=0$. Thus diverging widths and bounded endpoint motions preclude nonextinct eventual uniformity, including rational cones with bounded periodic offsets. No positivity of every residue is needed.

\begin{theorem}[Sharper triangular certificate]
Set $L=\max(0,M)$, $R=\max(0,1-m)$, $v=\max(M,1-m)\ge1$. If $b\ge\max(q-1,L+R+q-2)$, then uniform rows $b,\ldots,b+\lceil q/v\rceil$ force row $b$ zero.
\end{theorem}
\textbf{Boundary proof.} Core targets give full residue convolution. At any other target the actual sum omits terms, so equality and nonnegativity force every omitted summand to vanish. A positive source residue and a supported $d>0$ therefore propagate their least representative by $d$ without wrapping modulo $q$: a wrap would create a positive missing predecessor below zero. This is impossible after $\lceil q/d\rceil$ transitions. For $d\le0$, the least distance from the moving top increases by $1-d$ without wrapping. Choose an extreme attaining $v$.

The certificate uses at most $q+1$ rows. It is sharp for \emph{prescribed later initial rows}: choose the single jump $1-v$, large $b$ divisible by $q$, and initial value one on residue zero. Exactly $\lceil q/v\rceil$ nonzero uniform rows precede failure. This is not origin-start sharpness. For origin kernels supported in $\{0,1\}$, binomial comparisons instead show no uniform row at $n\ge q+1$, while equal weights pass at $n=q$.

\textbf{Inverse use and limits.} A nonzero block of $t+1$ uniform rows excludes every proposed finite kernel meeting its core threshold with $tv\ge q$, without reconstructing weights. Layer-dependent kernels in one common finite range still satisfy the mass proof (products replace powers of $K$). Without a common range, $f_n=1_{[1-n,n]}$ produces $n!$ uniform rows; the homogeneous infinite kernel $f\equiv1$ does too. Bounded widths also permit persistence. Signed weights and infinitely separated uniform rows remain open extensions here. Fixed-set killed-walk theory [3] is related background, not a novelty determination.
\newpage
\section*{3. When observations identify the kernel}
\textbf{Full labelled data.} Over a field, triangular rows through $N$ recover $f$ on $W_N$ if $p_{n-1}(0)$ and $p_{n-1}(n-1)$ are nonzero. Initially $f(0)=p_1(0)$, $f(1)=p_1(1)$; the next extremes follow from
\begin{align*}
f(n)&=\frac{p_n(n)-\sum_{z=1}^{n-1}p_{n-1}(z)f(n-z)}{p_{n-1}(0)},\\
f(1-n)&=\frac{p_n(0)-\sum_{z=0}^{n-2}p_{n-1}(z)f(-z)}{p_{n-1}(n-1)}.
\end{align*}
This uses full preceding rows, not only boundary sequences. Ordinary indicator counts with $0,1\in D$ give pivots; reduction to positive characteristic can cancel them.

\textbf{Funnel extension.} Assuming both preceding corner counts are nonzero, recovery works for arbitrary monotone widths, identifying exactly $[c-s_{N-1},c+s_N]$. After shearing, suppose the old window is $[-a,b]$ and the new one $[-b,k]$. Targets $b-a-1,\ldots,0$ recover new negative weights in descending order using the top pivot; targets $b+1,\ldots,k$ recover new positive weights in ascending order using the bottom pivot. Plateaus and sudden widening are permitted.

\textbf{A scalar inverse theorem.} On these funnels, for nonnegative real kernels, all actual rows through $N$ are constant and nonzero iff $f=\alpha>0$ on the observable window. Indeed each nonzero scalar row tests the entire one-dimensional constant-input space, so funnel rigidity applies. The common count is $\alpha^n\prod_{j=0}^{n-1}(s_j+1)$. Uniformity identifies kernel shape; one row sum fixes $\alpha$. Unbounded widths give a global converse.

\textbf{Compressed parity observations.} In the triangle let $f\ge0$ and $f(0)=f(1)=t>0$. Actual $2$-uniformity through $N$ forces $f=t$ on $W_N$ once $N\ge6$ for nonnegative real weights, $N\ge5$ for rationals, or $N\ge3$ for indicators. Each threshold is optimal. Normalize $t=1$ and put $x=f(-1)=f(2)$. The next forced pairs are
\[
c=\frac{x^2-x+2}{x+1},\qquad a=\frac{3x-1}{x+1}.
\]
A row-five interior comparison forces $(x-1)^2(x^2+2x-1)=0$. The extra real branch $x=\sqrt2-1$ gives a positive uniform prefix through five, but row six would require $f(-5)=7-5\sqrt2<0$. For indicators, already row three excludes $x=0$. Once a periodic seed is obtained, interior agreement with its periodic continuation, boundary uniformity and nonzero pivots propagate it indefinitely.

For indicators containing $0,1$, modulus three has sharp horizon four. Row-three equality forces $f(-1)=f(2)=x$, $f(-2)=f(3)=c$. With $P=1+3x+(1+x)c$, $Q=3+2x+x^2$, row four requires $f(-3)=f(4)=c+(Q/P)(1-c)$. Binary values force $c=1$, giving residue mask $(1,1,x)$. The same periodic-seed argument propagates it. General moduli and eventual-only uniformity remain open.

\textbf{Information lost.} Full-data kernel fibers are affine: $A(p)g=p$, with entries $A(p)_{(n,y),d}=p_{n-1}(y-d)$ when admitted. Row-space tests precisely describe field-valued identifiability; genuinely hidden intermediate rows yield polynomial fibers instead. Reflection $f(d)\mapsto f(1-d)$ reverses each row, so multisets and totals cannot distinguish $\{0,1,2\}$ from $\{-1,0,1\}$ despite positive reachability. Finite data cannot determine invisible tails or certify finite support. These delimit the inverse problem; general inverse-walk reconstruction is established research [4--6].
\newpage
\section*{4. Recognition, development and evidence}
\textbf{Recognition questions.}
\begin{enumerate}
\item Is the finite-range boundary certificate a familiar positive-operator or killed-walk lemma? Does its support-dependent horizon have useful content beyond standard arguments?
\item Are the growing-window/funnel converses and corner reconstruction recognizable consequences of structured-matrix or inverse-walk theory? Are the small-modulus compressed inverse tests worth isolating?
\end{enumerate}
\textbf{Development questions.}
\begin{enumerate}
\item What hypotheses force a periodic seed for general $q$? We have all-row indicator converses for $q=2,3$ with $0,1\in D$, but a positive nonperiodic $q=5$ prefix survives through six. Eventual uniformity is a separate question.
\item Can our constrained-walk category organize useful observation equivalences or quotients? Reflection is an explicit ambiguity; endpoint monoids alone lose generator information. Which geometry or observations restore useful identifiability?
\end{enumerate}

\textbf{What may already be known.} Equitable invariant spaces, linear observation fibers and recursive elimination are established methods. Reflection is an elementary walk symmetry. The potential content lies in the exact expanding-window converse, the support-dependent certificate and the compressed small-modulus tests. The inspected fixed-graph inverse theorems [4--6] use different observations and hypotheses; that distinction establishes neither originality nor independence from all their consequences. I would particularly welcome recognition of standard results that subsume these statements.

\textbf{Verification and AI assistance.} OpenAI Codex assisted with derivations, source comparison, exact computations, drafting and separate internal adversarial reviews. These are AI checks with shared limitations, not external peer review. The package includes finite exact tests, explicit-word comparisons and the exceptional branch over $\mathbb Q(\sqrt2)$; the universal claims rest on the displayed proofs. Source-reading boundaries and an assumption-by-assumption audit are supplied separately. Nicholas Ham remains the sole author; no participation by James is assumed.

\textbf{Selected primary references.}
\begin{enumerate}
\item East--Ham, \emph{Lattice paths and submonoids of $\Z^2$}, \href{https://arxiv.org/abs/1811.05735}{arXiv:1811.05735}: v1 Example 3.14 and \S4; v2/v3 Lemma 3.2 and Remark 3.4. The factorial examples are prior work.
\item Bailey et al., \href{https://arxiv.org/abs/1802.01001}{arXiv:1802.01001}, \S2: equitable/invariant-space background.
\item Lawler--Limic, \emph{Random Walk: A Modern Introduction}, \href{https://www.math.uchicago.edu/~lawler/srwbook.pdf}{author-hosted text}, \S5.2 pp.112--114 and \S12.4.1 pp.272--273: killed-walk definitions and fixed-set operators.
\item Liu--Siemons, \href{https://arxiv.org/abs/1911.00062v3}{arXiv:1911.00062v3}, Theorems 1.2, 5.1, 5.5 and 6.1: fixed-graph walk data/equivalence.
\item Bl{\aa}sten--Isozaki--Lassas--Lu, \href{https://arxiv.org/abs/2107.00494v2}{arXiv:2107.00494v2}, Condition 1.4, Theorems 1.5--1.8 and observation definitions, PDF pp.1--9: reversible finite-graph recovery under a Two-Points Condition.
\item De la Pe{\~n}a--Gzyl--McDonald, \href{https://arxiv.org/abs/math/0610821v1}{arXiv:math/0610821v1}, introduction and Theorem 1.1, PDF pp.1--3: augmented-tree recovery from hitting data.
\end{enumerate}
\end{document}
